% Options for packages loaded elsewhere
% Options for packages loaded elsewhere
\PassOptionsToPackage{unicode}{hyperref}
\PassOptionsToPackage{hyphens}{url}
\PassOptionsToPackage{dvipsnames,svgnames,x11names}{xcolor}
%
\documentclass[
  12pt,
]{article}
\usepackage{xcolor}
\usepackage[margin=1in]{geometry}
\usepackage{amsmath,amssymb}
\setcounter{secnumdepth}{-\maxdimen} % remove section numbering
\usepackage{iftex}
\ifPDFTeX
  \usepackage[T1]{fontenc}
  \usepackage[utf8]{inputenc}
  \usepackage{textcomp} % provide euro and other symbols
\else % if luatex or xetex
  \usepackage{unicode-math} % this also loads fontspec
  \defaultfontfeatures{Scale=MatchLowercase}
  \defaultfontfeatures[\rmfamily]{Ligatures=TeX,Scale=1}
\fi
\usepackage{lmodern}
\ifPDFTeX\else
  % xetex/luatex font selection
\fi
% Use upquote if available, for straight quotes in verbatim environments
\IfFileExists{upquote.sty}{\usepackage{upquote}}{}
\IfFileExists{microtype.sty}{% use microtype if available
  \usepackage[]{microtype}
  \UseMicrotypeSet[protrusion]{basicmath} % disable protrusion for tt fonts
}{}
\usepackage{setspace}
\makeatletter
\@ifundefined{KOMAClassName}{% if non-KOMA class
  \IfFileExists{parskip.sty}{%
    \usepackage{parskip}
  }{% else
    \setlength{\parindent}{0pt}
    \setlength{\parskip}{6pt plus 2pt minus 1pt}}
}{% if KOMA class
  \KOMAoptions{parskip=half}}
\makeatother
% Make \paragraph and \subparagraph free-standing
\makeatletter
\ifx\paragraph\undefined\else
  \let\oldparagraph\paragraph
  \renewcommand{\paragraph}{
    \@ifstar
      \xxxParagraphStar
      \xxxParagraphNoStar
  }
  \newcommand{\xxxParagraphStar}[1]{\oldparagraph*{#1}\mbox{}}
  \newcommand{\xxxParagraphNoStar}[1]{\oldparagraph{#1}\mbox{}}
\fi
\ifx\subparagraph\undefined\else
  \let\oldsubparagraph\subparagraph
  \renewcommand{\subparagraph}{
    \@ifstar
      \xxxSubParagraphStar
      \xxxSubParagraphNoStar
  }
  \newcommand{\xxxSubParagraphStar}[1]{\oldsubparagraph*{#1}\mbox{}}
  \newcommand{\xxxSubParagraphNoStar}[1]{\oldsubparagraph{#1}\mbox{}}
\fi
\makeatother


\usepackage{longtable,booktabs,array}
\newcounter{none} % for unnumbered tables
\usepackage{calc} % for calculating minipage widths
% Correct order of tables after \paragraph or \subparagraph
\usepackage{etoolbox}
\makeatletter
\patchcmd\longtable{\par}{\if@noskipsec\mbox{}\fi\par}{}{}
\makeatother
% Allow footnotes in longtable head/foot
\IfFileExists{footnotehyper.sty}{\usepackage{footnotehyper}}{\usepackage{footnote}}
\makesavenoteenv{longtable}
\usepackage{graphicx}
\makeatletter
\newsavebox\pandoc@box
\newcommand*\pandocbounded[1]{% scales image to fit in text height/width
  \sbox\pandoc@box{#1}%
  \Gscale@div\@tempa{\textheight}{\dimexpr\ht\pandoc@box+\dp\pandoc@box\relax}%
  \Gscale@div\@tempb{\linewidth}{\wd\pandoc@box}%
  \ifdim\@tempb\p@<\@tempa\p@\let\@tempa\@tempb\fi% select the smaller of both
  \ifdim\@tempa\p@<\p@\scalebox{\@tempa}{\usebox\pandoc@box}%
  \else\usebox{\pandoc@box}%
  \fi%
}
% Set default figure placement to htbp
\def\fps@figure{htbp}
\makeatother





\setlength{\emergencystretch}{3em} % prevent overfull lines

\providecommand{\tightlist}{%
  \setlength{\itemsep}{0pt}\setlength{\parskip}{0pt}}



 


\makeatletter
\@ifpackageloaded{caption}{}{\usepackage{caption}}
\AtBeginDocument{%
\ifdefined\contentsname
  \renewcommand*\contentsname{Table of contents}
\else
  \newcommand\contentsname{Table of contents}
\fi
\ifdefined\listfigurename
  \renewcommand*\listfigurename{List of Figures}
\else
  \newcommand\listfigurename{List of Figures}
\fi
\ifdefined\listtablename
  \renewcommand*\listtablename{List of Tables}
\else
  \newcommand\listtablename{List of Tables}
\fi
\ifdefined\figurename
  \renewcommand*\figurename{Figure}
\else
  \newcommand\figurename{Figure}
\fi
\ifdefined\tablename
  \renewcommand*\tablename{Table}
\else
  \newcommand\tablename{Table}
\fi
}
\@ifpackageloaded{float}{}{\usepackage{float}}
\floatstyle{ruled}
\@ifundefined{c@chapter}{\newfloat{codelisting}{h}{lop}}{\newfloat{codelisting}{h}{lop}[chapter]}
\floatname{codelisting}{Listing}
\newcommand*\listoflistings{\listof{codelisting}{List of Listings}}
\makeatother
\makeatletter
\makeatother
\makeatletter
\@ifpackageloaded{caption}{}{\usepackage{caption}}
\@ifpackageloaded{subcaption}{}{\usepackage{subcaption}}
\makeatother
\usepackage{bookmark}
\IfFileExists{xurl.sty}{\usepackage{xurl}}{} % add URL line breaks if available
\urlstyle{same}
\hypersetup{
  pdftitle={Problem Set 1},
  pdfauthor={Jay Ballesteros},
  colorlinks=true,
  linkcolor={blue},
  filecolor={Maroon},
  citecolor={Blue},
  urlcolor={Blue},
  pdfcreator={LaTeX via pandoc}}


\title{Problem Set 1}
\usepackage{etoolbox}
\makeatletter
\providecommand{\subtitle}[1]{% add subtitle to \maketitle
  \apptocmd{\@title}{\par {\large #1 \par}}{}{}
}
\makeatother
\subtitle{PPHA 34410 Corporate Finance}
\author{Jay Ballesteros}
\date{April 13, 2026}
\begin{document}
\maketitle


\setstretch{1.15}
\subsection{Part 1. Multiple Choice
Questions}\label{part-1.-multiple-choice-questions}

\textbf{1) Disney's new CEO emphasizes theme parks and streaming as core
growth drivers, while the company faces mixed investor sentiment and a
declining stock price. Which is the best interpretation of this
strategy?}

B. Disney is doubling down on its strongest proven segments to stabilize
investor expectations.

\textbf{2) SpaceX historically resisted going public but is now pursuing
an IPO to fund space-based data centers. What is the most accurate
explanation?)}

C. The company needs large, immediate capital for a high-risk,
capital-intensive project

\textbf{3) Which of the following represents potential benefits and
drawbacks of going public for SpaceX? Select all that apply}

A. Access to large-scale capital B. Increased short-term performance
pressure D. Greater liquidity for employees

\textbf{4) SpaceX conducts insider share sales at
\textasciitilde\$420/share to ``level-set'' valuation ahead of IPO. What
is the primary economic function of this?}

B. Providing a credible market-based signal of valuation

\textbf{5) Neom executives used overly optimistic assumptions and
shielded leadership from bad news. What is the core failure?}

C. Principal-agent problem with distorted information flow

\textbf{6) Which factors contributed to Neom's challenges? Select all
that apply.}

A. Unrealistic assumptions in financial models B. Information asymmetry
between leadership and operators D. Incentives to justify continued
investment

\textbf{7) Which best explains why SpaceX's ambitious projects may be
more credible than Neom's?}

B. SpaceX has demonstrated execution capability and cash flow generation

\textbf{8) Which conditions increase the probability that large,
visionary projects succeed? Select all that apply.}

A. Strong feedback loops B. Transparent information flow D. Proven
intermediate milestones

\begin{center}\rule{0.5\linewidth}{0.5pt}\end{center}

\subsection{Part 2. Short Calculation Questions (show your work for full
credit)}\label{part-2.-short-calculation-questions-show-your-work-for-full-credit}

\textbf{9) Calculate how much money (future value) you will have if you
invest \$100,000 today. (Hint: think about ``the rule of 72''.)}

\begin{enumerate}
\def\labelenumi{\alph{enumi}.}
\tightlist
\item
  at 10\% annually for 7.2 years: 198,621
\item
  at 4.5\% annually for 16 years: 202,237
\item
  at 8.0\% annually for 9 years: 199,900
\item
  at 9.0\% annually for 8 years: 199,256
\end{enumerate}

\textbf{10) Calculate the future value of \$100,000 invested today,
given the following rates and time periods, compounding annually}

\begin{enumerate}
\def\labelenumi{\alph{enumi}.}
\tightlist
\item
  r = 4\%, t = 5 years: 121,665
\item
  r = 10\%, t = 5 years: 161,051
\item
  r = 4\%, t = 20 years: 219,110
\item
  r = 10\%, t = 20 years: 672,749
\end{enumerate}

\textbf{11) Calculate the present value of \$100,000 received in t
number of years, given the following rates and time periods, compounding
annually}

\begin{enumerate}
\def\labelenumi{\alph{enumi}.}
\tightlist
\item
  r = 4\%, t = 5 years: 82,192.90
\item
  r = 10\%, t = 5 years: 62,092.13
\item
  r = 4\%, t = 20 years: 45,639.17
\item
  r = 10\%, t = 20 years: 14,864.36
\end{enumerate}

\textbf{12) You are building a new office headquarters that is estimated
to create efficiencies, saving your company \$25,000,000 annually for
ten years. (Assume the first savings appear one year from completion.)
You will move in immediately upon completion. The building cost, which
you will pay upon completion, is \$140,000,000. If your company's cost
of capital is 11\% annually, does this appear to be a good investment?
Why or why not?}

Considering that the present value of the savings is 147,232,250, then
the NPV is \(-140,000,000 + 147,232,250 = 7,232,250\). Since the NPV is
positive, this is a good investment.

\textbf{13) Perpetuity and Annuity}

\begin{enumerate}
\def\labelenumi{\alph{enumi}.}
\tightlist
\item
  You have the opportunity to invest in a growing perpetuity paying
  \$100,000 per annum in year 1. The rate of growth is expected to be
  6\% annually. If the market rate of interest is 9\% annually, how much
  should you be willing to pay for this perpetuity?
\end{enumerate}

Considering the information given, the present is
\(100,000/(0.09-0.06) = 3,333,333.33\). Therefore I will be willing to
pay 3,333,333.33 for this perpetuity.

\begin{enumerate}
\def\labelenumi{\alph{enumi}.}
\setcounter{enumi}{1}
\tightlist
\item
  You have the opportunity to invest in a growing 10 year annuity paying
  \$100,000 per annum in year 1. The growth rate is expected to be 6\%
  annually. If the market rate of interest is 9\% annually, how much
  should you be willing to pay for this annuity? What would you pay for
  the annuity if instead of 10 years, the annuity's term is 20 years?
\end{enumerate}

For the 10 year annuity, the present value is
\(100,000 * (1 - (1.06/1.09)^{10}) / (0.09 - 0.06) = 811,653\). For the
20 year annuity, the present value is
\(100,000 * (1 - (1.06/1.09)^{20}) / (0.09 - 0.06) = 1,425,673\).

\textbf{14) You won your state lottery and the prize is available in
several alternative payments. Which would you choose if the available
market interest rate is 6\% annually? Which, if the rate were 10\%? Show
your calculations.}

\begin{enumerate}
\def\labelenumi{\alph{enumi}.}
\tightlist
\item
  \$100,000,000 today: At 6\%: \$100,000,000 At 10\%: \$100,000,000
\item
  \$18,000,000 per annum for 8 years: At 6\%:
  \(18,000,000 * (1 - (1.06)^{-8}) / 0.06 = 111,776,220\) At 10\%:
  \(18,000,000 * (1 - (1.10)^{-8}) / 0.10 = 96,028,740\)
\item
  \$6,000,000 per annum in perpetuity: At 6\%:
  \(6,000,000 / 0.06 = 100,000,000\) At 10\%:
  \(6,000,000 / 0.10 = 60,000,000\)
\item
  \$10,000,000 per annum for 20 years: At 6\%:
  \(10,000,000 * (1 - (1.06)^{-20}) / 0.06 = 114,699,200\) At 10\%:
  \(10,000,000 * (1 - (1.10)^{-20}) / 0.10 = 85,135,600\)
\item
  \$2,500,000 in year one, growing at 3.5\% per annum in perpetuity At
  6\%: \(2,500,000 / (0.06 - 0.035) = 100,000,000\) At 10\%:
  \(2,500,000 / (0.10 - 0.035) = 38,461,538\)
\end{enumerate}

At 6\% interest rate, I will choose option d (\$114,699,200), as it has
the highest PV. At 10\%, option a (\$100,000,000 today), since higher
discount rates reduce the value of future payments significantly.

\textbf{15) Robert Downey Jr.~is reportedly earning over \$100 million
to return as Doctor Doom in Avengers: Doomsday and Avengers: Secret
Wars. Instead of taking the \$100 million in a lump sum payment in 2026,
Downey is considering an option where he would earn \$8 million per year
for 25 years starting in 2035. What is the internal rate of return he
would earn by forfeiting the immediate \$100 million in favor of the
deferred stream of income?}

We need to find r such that
\(100,000,000 = [8,000,000 * (1 - (1+r)^{-25}) / r] / (1+r)^8\). By
trial and error, at r = 3.5\%, the right side equals approximately
\$100,137,000, which is very close to 100,000,000. Therefore, the IRR is
approximately 3.5\%. This is a low return, suggesting Downey would be
better off taking the \$100 million lump sum today.




\end{document}
